\section{Introduction}

\subsection{Background and Motivation}
In modern data science and machine learning workflows, empirical studies consistently reveal that data practitioners spend between 60\% to 80\% of their total project lifecycle on exploratory data analysis (EDA), data cleaning, and feature preprocessing \cite{mckinney2010}. The fundamental premise of statistical learning, frequently encapsulated as ``Garbage In, Garbage Out'' (GIGO), dictates that even state-of-the-art predictive algorithms will yield degraded performance, high variance, or biased predictions if fed corrupted, noisy, or unstandardized inputs.

Real-world tabular datasets—originating from heterogeneous sources such as transactional enterprise databases, web scrapers, IoT loggers, and legacy spreadsheets—are plagued with systemic anomalies:
\begin{enumerate}[label=\arabic*.]
  \item \textbf{Missing and Null Inconsistencies:} Databases frequently employ divergent representations for missingness, including literal strings such as \texttt{"N/A"}, \texttt{"null"}, \texttt{"?"}, and whitespace dashes, preventing standard numerical routines from parsing columns reliably.
  \item \textbf{Formatting and Currency Contamination:} Quantitative columns containing financial transactions (e.g., \texttt{"\$1,200.50"}, \texttt{"1,500 đ"}, \texttt{"2500 VND"}) are imported as generic textual objects rather than floating-point numeric types.
  \item \textbf{Outliers and Distributional Skewness:} Extreme observations and recording errors severely distort sample variance, sample means, and normal distributions required by parametric models.
  \item \textbf{Multicollinearity and Redundancy:} High correlation between predictive features ($|r| \ge 0.85$) destabilizes matrix inversion in linear estimators and inflates standard errors.
  \item \textbf{Unencoded Categorical Data:} Non-numerical variables cannot be directly supplied to tensor operations without robust categorical encoding, which itself risks the dummy variable trap.
\end{enumerate}

Traditional workarounds rely heavily on ad-hoc Jupyter Notebook scripts written in Python or R. While flexible, this approach demands repetitive boilerplate code, lacks systematic reproducibility, and presents steep barriers to non-programming domain experts. Conversely, commercial solutions and emerging cloud-based Generative AI tools introduce severe data confidentiality vulnerabilities, high latency, non-deterministic transformations (hallucinations), and vendor lock-in.

To address these critical challenges, our group conceived and developed \textbf{DataSanitizer}—a high-performance, browser-accessible, deterministic data sanitization and profiling studio. By unifying an asynchronous FastAPI back-end with a reactive vanilla JavaScript Single-Page Application (SPA), DataSanitizer provides intuitive automated data health diagnosis and an interactive, ordered transformation pipeline engine.

\subsection{Objectives}
The primary objective of the DataSanitizer project is to deliver a reliable, lightweight, and auditable web platform for automated tabular data preprocessing. Specifically, the project establishes four measurable technical milestones:
\begin{itemize}
  \item \textbf{Objective 1 (Automated Health Profiling):} Construct an automated dataset profiling service capable of ingesting tabular files (\texttt{.csv} and \texttt{.xlsx} up to 10~MB) and computing comprehensive statistical diagnostics—including sorted null-percentage distributions, duplicated row counts, unique cardinality, and pairwise multicollinearity warnings ($|r| \ge 0.85$)—within 2.0 seconds.
  \item \textbf{Objective 2 (Ordered Enterprise Pipeline Execution):} Implement an in-memory transformation engine based on Pandas \cite{mckinney2010} and Scikit-learn \cite{pedregosa2011} that executes user-composed, ordered preprocessing pipelines featuring missing value imputation, robust outlier handling (IQR and Z-score winsorization/removal), categorical encoding (Label and One-Hot with collinearity avoidance), feature scaling, currency extraction, and floating-point rounding.
  \item \textbf{Objective 3 (Defensive Client-Side Validation):} Design a reactive user interface featuring context-aware dynamic type guarding that automatically inspects data types and prevents users from executing mathematically impermissible operations (e.g., executing numeric scaling on text columns), eliminating unhandled HTTP 400 server rejections.
  \item \textbf{Objective 4 (Transparent Auditability and Export):} Establish a verifiable audit logging mechanism that records pre- and post-processing row and column differentials alongside step-by-step transformation receipts, facilitating clean dataset downloads in standard CSV and JSON formats.
\end{itemize}

\subsection{Scope and Target Users}
\textbf{Target Users:}
DataSanitizer is engineered for data analysts, academic researchers, machine learning practitioners, and business intelligence specialists who require rapid, deterministic data preparation without constructing bespoke preprocessing scripts for every incoming dataset.

\textbf{In-Scope Features:}
\begin{itemize}
  \item Multipart file ingestion supporting CSV and modern Excel (XLSX) structures with strict size and MIME verification.
  \item Full-table health inspection and lazy-loaded single-column deep statistical inspection (mean, median, standard deviation, quartiles, IQR outlier count).
  \item Interactive construction, reordering, and removal of transformation pipeline steps.
  \item Robust currency string stripping and numerical casting via compiled regular expressions.
  \item Automated transformation audit summary generation and clean file export.
\end{itemize}

\textbf{Out-of-Scope Features:}
As an educational and departmental prototype, DataSanitizer currently excludes multi-tenant enterprise authentication (OAuth2/SSO), persistent relational database storage, long-term cloud object retention, and distributed cluster computation (e.g., Apache Spark).

\subsection{Technologies}
The technology stack for DataSanitizer was selected based on strict requirements for high execution speed, mathematical precision, standards compliance, and lightweight deployment footprints without external framework bloat. Table~\ref{tab:tech} summarizes the core technologies utilized.

\begin{table}[ht]
\centering\small
\renewcommand{\arraystretch}{1.2}
\begin{tabular}{@{}lll@{}}
\toprule
\textbf{Layer / Component} & \textbf{Technology} & \textbf{Technical Purpose}\\
\midrule
\textbf{Front-end UI}       & HTML5, Modern CSS3      & Structural semantics, modern cards, and responsive grid layout\\
\textbf{Front-end Logic}    & Vanilla JavaScript (ES6+) & Modular architecture, Fetch API, DOM manipulation, type guarding\\
\textbf{Back-end Framework} & FastAPI (Python 3.10+)  & Asynchronous ASGI routing, OpenAPI documentation, CORS\\
\textbf{Data Validation}    & Pydantic v2             & Strict typed schema parsing, HTTP 422 payload protection\\
\textbf{Analytical Core}    & Pandas \& NumPy         & Vectorized matrix manipulation, in-memory DataFrame transforms\\
\textbf{Machine Learning}   & Scikit-learn            & Standard/MinMax feature scalers, Label and One-Hot encoders\\
\textbf{Spreadsheet Parser} & OpenPyXL \& Zipfile     & Excel file decompression, structure and MIME verification\\
\textbf{Automated Testing}  & Pytest \& HTTPX         & Contract testing, pipeline verification, edge-case coverage\\
\textbf{Version Control}    & Git \& GitHub           & Collaborative branching, issue tracking, and code reviews\\
\bottomrule
\end{tabular}
\caption{Technologies utilized across the DataSanitizer platform.}
\label{tab:tech}
\end{table}

\subsection{Report Structure}
The remainder of this report is organized as follows:
\begin{itemize}
  \item \textbf{Section~2 (Analysis and Design):} Details functional and non-functional requirements, UI/UX interaction workflows, multi-tiered system architecture, data storage strategy, and RESTful API specifications.
  \item \textbf{Section~3 (Implementation):} Explains project modularity, key client-side type-guarding mechanisms, server-side data processing algorithms, and teamwork workflows.
  \item \textbf{Section~4 (Results and Conclusion):} Demonstrates automated test suite validation, evaluates objective fulfillment, discusses known constraints, and outlines future architectural improvements.
  \item \textbf{Appendix:} Provides comprehensive local execution instructions, API payload contracts, and ethical declarations regarding AI assistance.
\end{itemize}
